% Name: Huriye Kayadelen; Student number: 2200005795
% Options for packages loaded elsewhere
\PassOptionsToPackage{unicode}{hyperref}
\PassOptionsToPackage{hyphens}{url}
\documentclass[
  10pt,
  a4paper,
]{article}
\usepackage{xcolor}
\usepackage[margin=22mm]{geometry}
\usepackage{amsmath,amssymb}
\setcounter{secnumdepth}{-\maxdimen} % remove section numbering
\usepackage{iftex}
\ifPDFTeX
  \usepackage[T1]{fontenc}
  \usepackage[utf8]{inputenc}
  \usepackage{textcomp} % provide euro and other symbols
\else % if luatex or xetex
  \usepackage{unicode-math} % this also loads fontspec
  \defaultfontfeatures{Scale=MatchLowercase}
  \defaultfontfeatures[\rmfamily]{Ligatures=TeX,Scale=1}
\fi
\usepackage{lmodern}
\ifPDFTeX\else
  % xetex/luatex font selection
\fi
% Use upquote if available, for straight quotes in verbatim environments
\IfFileExists{upquote.sty}{\usepackage{upquote}}{}
\IfFileExists{microtype.sty}{% use microtype if available
  \usepackage[]{microtype}
  \UseMicrotypeSet[protrusion]{basicmath} % disable protrusion for tt fonts
}{}
\makeatletter
\@ifundefined{KOMAClassName}{% if non-KOMA class
  \IfFileExists{parskip.sty}{%
    \usepackage{parskip}
  }{% else
    \setlength{\parindent}{0pt}
    \setlength{\parskip}{6pt plus 2pt minus 1pt}}
}{% if KOMA class
  \KOMAoptions{parskip=half}}
\makeatother
\usepackage{listings}
\newcommand{\passthrough}[1]{#1}
\lstset{defaultdialect=[5.3]Lua}
\lstset{defaultdialect=[x86masm]Assembler}
\usepackage{longtable,booktabs,array}
\usepackage{caption}
\captionsetup[table]{skip=6pt}
\newcounter{none} % for unnumbered tables
\usepackage{calc} % for calculating minipage widths
% Correct order of tables after \paragraph or \subparagraph
\usepackage{etoolbox}
\makeatletter
\patchcmd\longtable{\par}{\if@noskipsec\mbox{}\fi\par}{}{}
\makeatother
\setlength{\emergencystretch}{3em} % prevent overfull lines
\providecommand{\tightlist}{%
  \setlength{\itemsep}{0pt}\setlength{\parskip}{0pt}}
% Name: Huriye Kayadelen; Student number: 2200005795
\usepackage{xcolor}
\usepackage{listings}
\usepackage{fancyhdr}
\usepackage{microtype}
\definecolor{garden}{HTML}{245A43}
\definecolor{codebg}{HTML}{F4F6F5}
\lstdefinelanguage{JavaScript}{keywords={const,let,function,return,if,else,for,of,in,new,try,catch,finally,throw,import,from,export,default,type,extends,infer,readonly,async,await,yield},sensitive=true,morecomment=[l]{//},morecomment=[s]{/*}{*/},morestring=[b]',morestring=[b]",morestring=[b]`}
\lstalias{js}{JavaScript}
\lstalias{ts}{JavaScript}
\lstalias{tsx}{JavaScript}
\lstalias{jsx}{JavaScript}
\lstalias{typescript}{JavaScript}
\lstalias{TypeScript}{JavaScript}
\lstdefinelanguage{text}{}
\lstdefinelanguage{powershell}{}
\lstdefinelanguage{sh}{}
\lstdefinelanguage{json}{}
\lstset{basicstyle=\ttfamily\scriptsize,keywordstyle=\color{garden}\bfseries,commentstyle=\color{gray},backgroundcolor=\color{codebg},frame=single,rulecolor=\color{black!12},breaklines=true,breakatwhitespace=false,columns=fullflexible,keepspaces=true,showstringspaces=false,tabsize=2,xleftmargin=3pt,xrightmargin=3pt,aboveskip=9pt,belowskip=9pt}
\AtBeginDocument{\hypersetup{colorlinks=true,linkcolor=garden,urlcolor=garden}}
\pagestyle{fancy}
\fancyhf{}
\fancyhead[L]{\small Destructuring Field Guide}
\fancyhead[R]{\small Huriye Kayadelen}
\fancyfoot[C]{\thepage}
\setlength{\headheight}{14pt}
\setlength{\emergencystretch}{3em}
\setcounter{tocdepth}{2}

% Keep each short example together across page boundaries.
\BeforeBeginEnvironment{lstlisting}{\par\begin{minipage}{\linewidth}}
\AfterEndEnvironment{lstlisting}{\end{minipage}\par}
\usepackage{bookmark}
\IfFileExists{xurl.sty}{\usepackage{xurl}}{} % add URL line breaks if available
\urlstyle{same}
% fallback for those not using the hyperref driver hyperxmp:
\makeatletter
\@ifundefined{xmpquote}{\newcommand{\xmpquote}[1]{#1}}{}
\makeatother
\hypersetup{
  pdftitle={Destructuring Field Guide},
  pdfauthor={Huriye Kayadelen (2200005795)},
  hidelinks,
  pdfcreator={LaTeX via pandoc}}

\title{Destructuring Field Guide}
\author{Huriye Kayadelen (2200005795)}
\date{16 November 2026}

\begin{document}
\maketitle

{
\setcounter{tocdepth}{3}
\tableofcontents
}
Name: Huriye Kayadelen Student number: 2200005795

Destructuring Field Guide

This guide covers all twelve JavaScript and strict TypeScript forms,
five pitfalls, T1-T5, React P1-P7 and the four bonus sections. The
examples use a balcony garden theme.

Bonuses submitted for assessment: 1 (LaTeX), 2 (Depth), 3 (Real code),
and 4 (Git repository).

\section{Run and scope}\label{run-and-scope}

Node.js 20.19 or newer is required.

\begin{lstlisting}[language=sh]
npm ci
npm run all
npm run check
\end{lstlisting}

strict and noUncheckedIndexedAccess stay enabled, and library checking
is enabled. The single permitted T5 line uses @ts-expect-error. The
deliberate compiler experiments are run separately without suppressions
and their temporary sources are removed. output.txt records the actual
run; logs/check.txt records the strict check.

\section{Field guide}\label{field-guide}

JS files contain the runtime examples; corresponding TS files add tuple,
object, parameter and collection types. The repeated form lines in
output.txt are the real JS and TS runs. Snippets below can be run
separately.

\subsection{Form 1: Values by position}\label{form-1-values-by-position}

\textbf{Syntax:}
\texttt{const [free,\allowbreak{} used] =\allowbreak{} traySpace(3,\allowbreak{} 4,\allowbreak{} 7)}

\textbf{Example} (js/arrays.mjs; focused snippet):

\begin{lstlisting}
function traySpace(rows, columns, planted) { return [rows * columns - planted, planted] }
const [free, used] = traySpace(3, 4, 7)
console.log('form 1:', free, used)
\end{lstlisting}

\textbf{Output:}

\begin{lstlisting}
form 1: 5 7
\end{lstlisting}

\textbf{Use it when:} a tray calculation returns two related values.

\textbf{Watch out:} Position decides the name. Reversing the bindings
reverses their meaning.

Plain assignment equivalent (in its own scope):

\begin{lstlisting}
const result = traySpace(3, 4, 7)
const free = result[0]
const used = result[1]
\end{lstlisting}

\subsection{Form 2: Skip and array
rest}\label{form-2-skip-and-array-rest}

\textbf{Syntax:}
\texttt{const [,\allowbreak{} next,\allowbreak{} .\allowbreak{}.\allowbreak{}.\allowbreak{}later] =\allowbreak{} sowingOrder}

\textbf{Example} (js/arrays.mjs; focused snippet):

\begin{lstlisting}
const sowingOrder = ['basil', 'mint', 'thyme', 'sage']
const [, next, ...later] = sowingOrder
console.log('form 2:', next, JSON.stringify(later))
\end{lstlisting}

\textbf{Output:}

\begin{lstlisting}
form 2: mint ["thyme","sage"]
\end{lstlisting}

\textbf{Use it when:} the first sowing is done and you need the
remaining queue.

\textbf{Watch out:} Rest must come last. It collects a new array, but
nested values would still be shared.

Plain assignment equivalent (in its own scope):

\begin{lstlisting}
const next = sowingOrder[1]
const later = sowingOrder.slice(2)
\end{lstlisting}

\subsection{Form 3: Swap}\label{form-3-swap}

\textbf{Syntax:}
\texttt{[morning,\allowbreak{} evening] =\allowbreak{} [evening,\allowbreak{} morning]}

\textbf{Example} (js/arrays.mjs; focused snippet):

\begin{lstlisting}
let morning = 'east', evening = 'west'
;[morning, evening] = [evening, morning]
console.log('form 3:', morning, evening)
\end{lstlisting}

\textbf{Output:}

\begin{lstlisting}
form 3: west east
\end{lstlisting}

\textbf{Use it when:} two pot positions exchange roles.

\textbf{Watch out:} Use let. The leading semicolon prevents this array
from joining the previous statement.

Plain assignment equivalent (in its own scope):

\begin{lstlisting}
const oldMorning = morning
morning = evening
evening = oldMorning
\end{lstlisting}

\subsection{Form 4: Array defaults}\label{form-4-array-defaults}

\textbf{Syntax:}
\texttt{const [missing =\allowbreak{} 6] =\allowbreak{} []}

\textbf{Example} (js/arrays.mjs; focused snippet):

\begin{lstlisting}
const [missing = 6] = []
const [unset = 6, explicitNull = 6] = [undefined, null]
console.log('form 4:', missing, unset, explicitNull)
\end{lstlisting}

\textbf{Output:}

\begin{lstlisting}
form 4: 6 6 null
\end{lstlisting}

\textbf{Use it when:} a missing seed count needs a fallback.

\textbf{Watch out:} Only undefined triggers the default. Null is kept.

Plain assignment equivalent (in its own scope):

\begin{lstlisting}
const empty = []
const values = [undefined, null]
const missing = empty[0] === undefined ? 6 : empty[0]
const unset = values[0] === undefined ? 6 : values[0]
const explicitNull = values[1] === undefined ? 6 : values[1]
\end{lstlisting}

\subsection{Form 5: Properties by name}\label{form-5-properties-by-name}

\textbf{Syntax:}
\texttt{const \{ crop,\allowbreak{} liters \} =\allowbreak{} bed}

\textbf{Example} (js/objects.mjs; focused snippet):

\begin{lstlisting}
const bed = { crop: 'basil', liters: 2 }
const { crop, liters } = bed
console.log('form 5:', crop, liters)
\end{lstlisting}

\textbf{Output:}

\begin{lstlisting}
form 5: basil 2
\end{lstlisting}

\textbf{Use it when:} a watering record has named fields.

\textbf{Watch out:} Property order does not matter. A missing property
gives undefined.

Plain assignment equivalent (in its own scope):

\begin{lstlisting}
const crop = bed.crop
const liters = bed.liters
\end{lstlisting}

\subsection{Form 6: Rename}\label{form-6-rename}

\textbf{Syntax:}
\texttt{const \{ crop:\allowbreak{} label \} =\allowbreak{} bed}

\textbf{Example} (js/objects.mjs; focused snippet):

\begin{lstlisting}
const bed = { crop: 'basil' }
const { crop: label } = bed
console.log('form 6:', label)
\end{lstlisting}

\textbf{Output:}

\begin{lstlisting}
form 6: basil
\end{lstlisting}

\textbf{Use it when:} a crop name is used as a display label.

\textbf{Watch out:} The binding is label. This pattern does not create a
crop variable.

Plain assignment equivalent (in its own scope):

\begin{lstlisting}
const label = bed.crop
\end{lstlisting}

\subsection{Form 7: Object defaults}\label{form-7-object-defaults}

\textbf{Syntax:}
\texttt{const \{ liters:\allowbreak{} missing =\allowbreak{} 1 \} =\allowbreak{} \{\}}

\textbf{Example} (js/objects.mjs; focused snippet):

\begin{lstlisting}
const { liters: missing = 1 } = {}
const { liters: unset = 1 } = { liters: undefined }
const { liters: explicitNull = 1 } = { liters: null }
console.log('form 7:', missing, unset, explicitNull)
\end{lstlisting}

\textbf{Output:}

\begin{lstlisting}
form 7: 1 1 null
\end{lstlisting}

\textbf{Use it when:} an omitted watering amount should use one liter.

\textbf{Watch out:} Missing and undefined use one; null does not. The
whole source must still be an object.

Plain assignment equivalent (in its own scope):

\begin{lstlisting}
const absent = {}
const undefinedField = { liters: undefined }
const nullField = { liters: null }
const missing = absent.liters === undefined ? 1 : absent.liters
const unset = undefinedField.liters === undefined ? 1 : undefinedField.liters
const explicitNull = nullField.liters === undefined ? 1 : nullField.liters
\end{lstlisting}

\subsection{Form 8: Object rest}\label{form-8-object-rest}

\textbf{Syntax:}
\texttt{const \{ id,\allowbreak{} .\allowbreak{}.\allowbreak{}.\allowbreak{}details \} =\allowbreak{} bed}

\textbf{Example} (js/objects.mjs; shared runnable example):

\begin{lstlisting}
const bed = { id: 'B2', crop: 'basil', liters: 2, sensor: { moisture: 34 }, days: ['Mon', 'Thu'] }
const { crop, liters } = bed
console.log('form 5:', crop, liters)
const { crop: label } = bed
console.log('form 6:', label)
const { liters: missing = 1 } = {}
const { liters: unset = 1 } = { liters: undefined }
const { liters: explicitNull = 1 } = { liters: null }
console.log('form 7:', missing, unset, explicitNull)
const { id, ...details } = bed
console.log('form 8:', id, JSON.stringify(details))
const { sensor: { moisture }, days: [firstDay] } = bed
console.log('form 9:', moisture, firstDay)
const key = 'liters'
const { [key]: selected } = bed
console.log('form 10:', key, selected)
\end{lstlisting}

\textbf{Output:}

\begin{lstlisting}
form 8: B2 {"crop":"basil","liters":2,"sensor":{"moisture":34},"days":["Mon","Thu"]}
\end{lstlisting}

\textbf{Use it when:} a public record should leave out the bed ID.

\textbf{Watch out:} Rest copies remaining own enumerable properties,
shallowly.

Plain assignment equivalent (in its own scope):

\begin{lstlisting}
const id = bed.id
const details = Object.assign({}, bed)
delete details.id
\end{lstlisting}

\subsection{Form 9: Nested object and
array}\label{form-9-nested-object-and-array}

\textbf{Syntax:}
\texttt{const \{ sensor:\allowbreak{} \{ moisture \},\allowbreak{} days:\allowbreak{} [firstDay] \} =\allowbreak{} bed}

\textbf{Example} (js/objects.mjs; shared runnable example):

\begin{lstlisting}
const bed = { id: 'B2', crop: 'basil', liters: 2, sensor: { moisture: 34 }, days: ['Mon', 'Thu'] }
const { crop, liters } = bed
console.log('form 5:', crop, liters)
const { crop: label } = bed
console.log('form 6:', label)
const { liters: missing = 1 } = {}
const { liters: unset = 1 } = { liters: undefined }
const { liters: explicitNull = 1 } = { liters: null }
console.log('form 7:', missing, unset, explicitNull)
const { id, ...details } = bed
console.log('form 8:', id, JSON.stringify(details))
const { sensor: { moisture }, days: [firstDay] } = bed
console.log('form 9:', moisture, firstDay)
const key = 'liters'
const { [key]: selected } = bed
console.log('form 10:', key, selected)
\end{lstlisting}

\textbf{Output:}

\begin{lstlisting}
form 9: 34 Mon
\end{lstlisting}

\textbf{Use it when:} a sensor reading and the first watering day belong
together.

\textbf{Watch out:} Missing sensor throws. An empty days array gives
undefined; keep deep patterns short.

Plain assignment equivalent (in its own scope):

\begin{lstlisting}
const moisture = bed.sensor.moisture
const firstDay = bed.days[0]
\end{lstlisting}

\subsection{Form 10: Computed key}\label{form-10-computed-key}

\textbf{Syntax:}
\texttt{const \{ [key]:\allowbreak{} selected \} =\allowbreak{} bed}

\textbf{Example} (js/objects.mjs; shared runnable example):

\begin{lstlisting}
const bed = { id: 'B2', crop: 'basil', liters: 2, sensor: { moisture: 34 }, days: ['Mon', 'Thu'] }
const { crop, liters } = bed
console.log('form 5:', crop, liters)
const { crop: label } = bed
console.log('form 6:', label)
const { liters: missing = 1 } = {}
const { liters: unset = 1 } = { liters: undefined }
const { liters: explicitNull = 1 } = { liters: null }
console.log('form 7:', missing, unset, explicitNull)
const { id, ...details } = bed
console.log('form 8:', id, JSON.stringify(details))
const { sensor: { moisture }, days: [firstDay] } = bed
console.log('form 9:', moisture, firstDay)
const key = 'liters'
const { [key]: selected } = bed
console.log('form 10:', key, selected)
\end{lstlisting}

\textbf{Output:}

\begin{lstlisting}
form 10: liters 2
\end{lstlisting}

\textbf{Use it when:} the selected metric is chosen by a key variable.

\textbf{Watch out:} The brackets evaluate key. In TS, keyof limits it to
valid property names.

Plain assignment equivalent (in its own scope):

\begin{lstlisting}
const selected = bed[key]
\end{lstlisting}

\subsection{Form 11: Object parameter
default}\label{form-11-object-parameter-default}

\textbf{Syntax:}
\texttt{function watering(\{ liters =\allowbreak{} 1,\allowbreak{} crop =\allowbreak{} 'seedlings' \} =\allowbreak{} \{\})}

\textbf{Example} (js/params.mjs; shared runnable example):

\begin{lstlisting}
function watering({ liters = 1, crop = 'seedlings' } = {}) { return crop + ': ' + liters + ' L' }
console.log('form 11:', watering(), '|', watering({ crop: 'basil', liters: 2 }))
// Form 11b: array parameter requested by starter.
function area([rows, columns]) { return rows * columns }
console.log('form 11b:', area([3, 4]))
const weekly = { Mon: 2, Thu: 3 }
for (const [day, liters] of Object.entries(weekly)) console.log('form 12 object:', day, liters)
const pots = new Map([['B2', 4], ['C1', 2]])
for (const [bed, count] of pots) console.log('form 12 Map:', bed, count)
\end{lstlisting}

\textbf{Output:}

\begin{lstlisting}
form 11: seedlings: 1 L | basil: 2 L
form 11b: 12
\end{lstlisting}

\textbf{Use it when:} callers may omit the whole options object or
individual fields.

\textbf{Watch out:} The outer default handles omitted/undefined
arguments, not null. Field defaults are a separate step. The typed
version uses optional fields so \{\} is valid.

Plain assignment equivalent (in its own scope):

\begin{lstlisting}
function watering(options = {}) {
 const liters = options.liters === undefined ? 1 : options.liters
 const crop = options.crop === undefined ? 'seedlings' : options.crop
 return crop + ': ' + liters + ' L'
}
\end{lstlisting}

\subsection{Form 12: Entry loops}\label{form-12-entry-loops}

\textbf{Syntax:}
\texttt{for (const [day,\allowbreak{} liters] of Object.\allowbreak{}entries(weekly))}

\textbf{Example} (js/params.mjs; shared runnable example):

\begin{lstlisting}
function watering({ liters = 1, crop = 'seedlings' } = {}) { return crop + ': ' + liters + ' L' }
console.log('form 11:', watering(), '|', watering({ crop: 'basil', liters: 2 }))
// Form 11b: array parameter requested by starter.
function area([rows, columns]) { return rows * columns }
console.log('form 11b:', area([3, 4]))
const weekly = { Mon: 2, Thu: 3 }
for (const [day, liters] of Object.entries(weekly)) console.log('form 12 object:', day, liters)
const pots = new Map([['B2', 4], ['C1', 2]])
for (const [bed, count] of pots) console.log('form 12 Map:', bed, count)
\end{lstlisting}

\textbf{Output:}

\begin{lstlisting}
form 12 object: Mon 2
form 12 object: Thu 3
form 12 Map: B2 4
form 12 Map: C1 2
\end{lstlisting}

\textbf{Use it when:} each watering day or Map entry needs a key and a
value.

\textbf{Watch out:} Object.entries gives string keys. A Map preserves
its key types and insertion order.

Plain assignment equivalent (in its own scope):

\begin{lstlisting}
for (const entry of Object.entries(weekly)) {
 const day = entry[0]
 const liters = entry[1]
 console.log(day, liters)
}
for (const entry of pots) {
 const bed = entry[0]
 const count = entry[1]
 console.log(bed, count)
}
\end{lstlisting}

\section{Rest and spread}\label{rest-and-spread}

Rest collects remaining items in a binding pattern:
\texttt{const \{ id,\allowbreak{} .\allowbreak{}.\allowbreak{}.\allowbreak{}details \} =\allowbreak{} bed}.
Spread expands values in an expression:
\texttt{const copy =\allowbreak{} \{ .\allowbreak{}.\allowbreak{}.\allowbreak{}bed \}}.
Both make shallow containers. Neither deep-copies sensor. Array rest
collects after a position; object rest excludes named properties.

\section{Five pitfalls}\label{five-pitfalls}

All five run in js/pitfalls.mjs. The thrown errors are deliberate
experiments and are caught.

\subsection{Pitfall 1}\label{pitfall-1}

\begin{lstlisting}
pitfall 1: undefined TypeError: Cannot destructure property 'liters' of 'source' as it is undefined.
pitfall 1: null TypeError: Cannot destructure property 'liters' of 'source' as it is null.
\end{lstlisting}

Neither undefined nor null can supply properties, so destructuring
throws TypeError.

\subsection{Pitfall 2}\label{pitfall-2}

\begin{lstlisting}
pitfall 2: undefined 1
pitfall 2: undefined 1
pitfall 2: null null
\end{lstlisting}

The default is used only for undefined; the first source is missing
liters, the second explicitly sets it to undefined, and null remains
null.

\subsection{Pitfall 3}\label{pitfall-3}

\begin{lstlisting}
pitfall 3: SyntaxError: Unexpected token '='
pitfall 3 fixed: 2
\end{lstlisting}

A leading brace is parsed as a block; parentheses turn the pattern into
an assignment expression. The invalid code is compiled inside new
Function so its SyntaxError can be caught.

\subsection{Pitfall 4}\label{pitfall-4}

\begin{lstlisting}
pitfall 4: 10 true {"crop":"basil"}
\end{lstlisting}

The sensor binding shares the original nested object; changing it
changes the original, while rest only creates a shallow outer object.

\subsection{Pitfall 5}\label{pitfall-5}

\begin{lstlisting}
pitfall 5: basil 2 | 2 basil
\end{lstlisting}

Object bindings follow names regardless of property order, while array
bindings follow positions.

\section{TypeScript T1-T5}\label{typescript-t1-t5}

Examples are in ts/types.ts. strict and noUncheckedIndexedAccess stay
enabled.

\subsection{T1}\label{t1}

In \{ a: string \}, string is the new variable name. The correct
annotation is const \{ a \}: \{ a: string \} = source. The two bindings
both contain basil.

\begin{lstlisting}
T1: basil basil
\end{lstlisting}

\subsection{T2}\label{t2}

tray(): {[}number, string{]} fixes the length and the type at each
position. as const infers readonly {[}4, "mint"{]}, keeping literal
types. Readonly is a compile-time restriction; it does not freeze the
runtime array.

\begin{lstlisting}
T2: 12 basil 4 mint
\end{lstlisting}

\subsection{T3}\label{t3}

An ordinary string{[}{]} may be empty. first is string \textbar{}
undefined. The branch checks undefined before toUpperCase. fallback is
string because its default removes undefined.

\begin{lstlisting}
T3: undefined empty no pot
\end{lstlisting}

\subsection{T4}\label{t4}

For label?: string \textbar{} null, the binding without a default has
type string \textbar{} null \textbar{} undefined. With a default it has
type string \textbar{} null. The actual null remains null; the missing
field becomes unlabelled.

\begin{lstlisting}
T4: null null unlabelled
\end{lstlisting}

\subsection{T5}\label{t5}

The deliberate experiment misspells liters as litres. Removing only the
suppression produced the actual diagnostic below. It was then restored
on that one line. The fixed binding uses liters. The runtime experiment
shows undefined for the wrong key and 2 for the corrected key. tsx runs
code but does not perform the separate type check.

\begin{lstlisting}
T5: undefined 2
\end{lstlisting}

Actual diagnostic, captured in logs/t5-diagnostic.txt (expected exit 2):

\begin{lstlisting}
.compiler-experiments/t5-unsuppressed.tsx(27,9): error TS2339: Property 'litres' does not exist on type '{ liters: number; }'.
\end{lstlisting}

\section{Props and children - Day 2}\label{props-and-children---day-2}

Both props files run P1-P7 using renderToStaticMarkup. The snippets
below are from ts/props.tsx; js/props.jsx uses the same components
without annotations. Shared types/imports:

\begin{lstlisting}
import type { ReactNode, PropsWithChildren, ComponentProps } from 'react'
import { renderToStaticMarkup } from 'react-dom/server'
type CardProps = { title: string; children?: ReactNode }
\end{lstlisting}

\subsection{P1}\label{p1}

\begin{lstlisting}
function CardA(props: CardProps) { return <section><h2>{props.title}</h2>{props.children}</section> }
function CardB({ title, children }: CardProps) { return <section><h2>{title}</h2>{children}</section> }
\end{lstlisting}

A component gets one props object. CardA reads its properties directly;
CardB binds the same properties in its parameter. That combines form 5
(named properties) with form 11 (parameter patterns). Neither version
changes what React receives.

Real output (the JS and TS runs agree):

\begin{lstlisting}
P1 A: <section><h2>Balcony</h2>basil</section>
P1 B: <section><h2>Balcony</h2>basil</section>
\end{lstlisting}

\subsection{P2}\label{p2}

\begin{lstlisting}
<CardB title="Balcony">mint</CardB>
<CardB title="Balcony" children="mint" />
\end{lstlisting}

The content between tags becomes the children property. These two calls
render the same HTML. Missing children is allowed and renders nothing;
it is not a second argument to the component.

Real output (the JS and TS runs agree):

\begin{lstlisting}
P2 content: <section><h2>Balcony</h2>mint</section>
P2 attribute: <section><h2>Balcony</h2>mint</section>
\end{lstlisting}

\subsection{P3}\label{p3}

\begin{lstlisting}
const cases: [string, ReactNode][] = [
  ['string', 'basil'], ['number', 3], ['element', <b>mint</b>],
  ['elements', [<b key="b">basil</b>, <i key="m">mint</i>]],
  ['undefined', undefined], ['null', null], ['false', false]
]
for (const [name, child] of cases) console.log('P3 ' + name + ':', renderToStaticMarkup(<section>{child}</section>))
function RenderCrop({ children }: { children: (crop: string) => ReactNode }) {
  return <section>{children('mint')}</section>
}
<RenderCrop>{crop => <b>{crop}</b>}</RenderCrop>
\end{lstlisting}

Strings and numbers render as text; elements render as HTML. Several
elements retain their order. undefined, null and false leave no visible
content. A function is not an ordinary ReactNode: passing it directly
warns and leaves the section empty. RenderCrop instead accepts a
function in its own explicit type and calls it with mint. The returned
element is what React renders. The direct function and object failures
are deliberate experiments in the JS helper invalidChildren(), also
called by the TS run, so TS needs no suppression or unsafe cast.

Real output (the JS and TS runs agree):

\begin{lstlisting}
P3 string: <section>basil</section>
P3 number: <section>3</section>
P3 element: <section><b>mint</b></section>
P3 elements: <section><b>basil</b><i>mint</i></section>
P3 undefined: <section></section>
P3 null: <section></section>
P3 false: <section></section>
P3 function: <section></section>
P3 object deliberate experiment: Error: Objects are not valid as a React child (found: object with keys {crop}). If you meant to render a collection of children, use an array instead.
P3 render prop: <section><b>mint</b></section>
\end{lstlisting}

The JS probe really passes
\texttt{\textless{}\allowbreak{}section\textgreater{}\allowbreak{}\{grow\}\textless{}\allowbreak{}/\allowbreak{}section\textgreater{}\allowbreak{}}
where
\texttt{grow =\allowbreak{} () =\allowbreak{}\textgreater{}\allowbreak{} \textless{}\allowbreak{}b\textgreater{}\allowbreak{}mint\textless{}\allowbreak{}/\allowbreak{}b\textgreater{}\allowbreak{}},
and
\texttt{\textless{}\allowbreak{}section\textgreater{}\allowbreak{}\{\{ crop:\allowbreak{} 'mint' \}\}\textless{}\allowbreak{}/\allowbreak{}section\textgreater{}\allowbreak{}}.
The console capture forwards the original React arguments to
console.log, so the raw \texttt{\%s} placeholders and following
substitution arguments are preserved:

\begin{lstlisting}
React warning: Functions are not valid as a React child. This may happen if you return %s instead of <%s /> from render. Or maybe you meant to call this function rather than return it. grow grow
\end{lstlisting}

\subsection{P4}\label{p4}

\begin{lstlisting}
function Frame({ title, children, ...rest }: CardProps & ComponentProps<'section'>) {
  return <section {...rest}><h2>{title}</h2>{children}</section>
}
const dom: ComponentProps<'section'> = { id: 'bed-B2', className: 'from-spread' }
<Frame title="Water" {...dom}>2 L</Frame>
<section {...dom} className="explicit">basil</section>
<section className="explicit" {...dom}>basil</section>
\end{lstlisting}

Rest collects all properties except title and children. Spreading rest
forwards id and className onto the section, while children remains body
content and title becomes a heading. React writes className as class in
HTML. When values conflict, the later property wins. The TS DOM type
makes className optional: it describes a normal forwarding component.

Real output (the JS and TS runs agree):

\begin{lstlisting}
P4 rest: <section id="bed-B2" class="from-spread"><h2>Water</h2>2 L</section>
P4 explicit last: <section id="bed-B2" class="explicit">basil</section>
P4 spread last: <section class="from-spread" id="bed-B2">basil</section>
\end{lstlisting}

\subsection{P5}\label{p5}

\begin{lstlisting}
function Slots({ header, footer, children }: PropsWithChildren<{ header: ReactNode; footer: ReactNode }>) {
  return <article><header>{header}</header><main>{children}</main><footer>{footer}</footer></article>
}
function Flow({ children }: PropsWithChildren) { return <article>{children}</article> }
<Slots header={<h2>Water</h2>} footer={<small>Evening</small>}><p>2 L</p></Slots>
<Flow><header><h2>Water</h2></header><main><p>2 L</p></main><footer><small>Evening</small></footer></Flow>
\end{lstlisting}

Slots fixes the header/body/footer structure and names the roles. I
would choose it for a consistent watering card whose callers supply
content but not layout. Flow leaves structure and order to the caller; I
would choose it for a flexible garden journal with mixed sections. They
can render identical HTML here, but Flow does not enforce the slot
layout.

Real output (the JS and TS runs agree):

\begin{lstlisting}
P5 slots: <article><header><h2>Water</h2></header><main><p>2 L</p></main><footer><small>Evening</small></footer></article>
P5 children: <article><header><h2>Water</h2></header><main><p>2 L</p></main><footer><small>Evening</small></footer></article>
\end{lstlisting}

\subsection{P6}\label{p6}

\begin{lstlisting}
function DefaultCard({ title = 'Balcony', children }: { title?: string | null; children?: ReactNode }) {
  return <section><h2>{title}</h2>{children}</section>
}
<DefaultCard />
<DefaultCard title={undefined} />
<DefaultCard title={null} />
\end{lstlisting}

This is the form 7 default inside a form 11 parameter. A missing title
and an explicit undefined both use Balcony. Null is an explicit value
and does not activate the default; React renders it as no heading
content. An empty string would also stay empty. A field default does not
protect destructuring an entire null props object.

Real output (the JS and TS runs agree):

\begin{lstlisting}
P6 missing: <section><h2>Balcony</h2></section>
P6 undefined: <section><h2>Balcony</h2></section>
P6 null: <section><h2></h2></section>
\end{lstlisting}

\subsection{P7}\label{p7}

\begin{lstlisting}
function GardenButton({ children, ...rest }: ComponentProps<'button'>) {
  return <button {...rest}>{children}</button>
}
<CardB title="Seeds">{3}</CardB>
<Flow />
<GardenButton type="button" disabled>Water</GardenButton>
\end{lstlisting}

children?: ReactNode permits missing children, text, numbers, elements,
arrays of nodes and empty values such as null/false. It rejects a
function or an arbitrary object. PropsWithChildren adds optional
ReactNode children to P; it does not make other required props optional
and does not turn a function into a node.
ComponentProps\textless\textquotesingle button\textquotesingle\textgreater{}
reuses React button props, including children, type, disabled and event
types. It rejects href (an anchor prop) and disabled="yes". Runtime
rendering alone does not check these types: tsx executes, while tsc
checks. The deliberate rejected examples were each compiled separately
in disposable files and removed; exact sources, exit codes and
diagnostics are in logs/compiler-diagnostics.txt.

Real output (the JS and TS runs agree):

\begin{lstlisting}
P7 ReactNode: <section><h2>Seeds</h2>3</section>
P7 PropsWithChildren: <article></article>
P7 button: <button type="button" disabled="">Water</button>
\end{lstlisting}

\section{JavaScript versus
TypeScript}\label{javascript-versus-typescript}

Types do not change destructuring at runtime. These are actual strict
compiler diagnostics, not predicted messages. Every case was an isolated
deliberate experiment. Full diagnostics and experiment sources:
\url{logs/compiler-diagnostics.txt}.

{\def\LTcaptype{none} % do not increment counter
\begin{longtable}[]{@{}
  >{\raggedright\arraybackslash}p{(\linewidth - 6\tabcolsep) * \real{0.1600}}
  >{\raggedright\arraybackslash}p{(\linewidth - 6\tabcolsep) * \real{0.2200}}
  >{\raggedright\arraybackslash}p{(\linewidth - 6\tabcolsep) * \real{0.2200}}
  >{\raggedright\arraybackslash}p{(\linewidth - 6\tabcolsep) * \real{0.4000}}@{}}
\toprule\noalign{}
\begin{minipage}[b]{\linewidth}\raggedright
Form
\end{minipage} & \begin{minipage}[b]{\linewidth}\raggedright
Same behavior
\end{minipage} & \begin{minipage}[b]{\linewidth}\raggedright
Type information
\end{minipage} & \begin{minipage}[b]{\linewidth}\raggedright
Observed compiler error
\end{minipage} \\
\midrule\noalign{}
\endhead
\bottomrule\noalign{}
\endlastfoot
1: tuple return & Positions keep their meaning & {[}number, string{]} &
TS2322: Type \textquotesingle number\textquotesingle{} is not assignable
to type \textquotesingle string\textquotesingle. \\
6 / T1: rename & The name after the colon is the binding & Type after
the whole pattern & TS2304: Cannot find name
\textquotesingle a\textquotesingle. \\
7: default & Only undefined activates the default & Include null when
allowed & TS18047: \textquotesingle label\textquotesingle{} is possibly
\textquotesingle null\textquotesingle. \\
9 / T3: array item & An empty array has no first item & string{[}{]}
gives string \textbar{} undefined & TS18048:
\textquotesingle first\textquotesingle{} is possibly
\textquotesingle undefined\textquotesingle. \\
10: computed key & The key expression selects a property & keyof typeof
bed & TS2322: Type \textquotesingle"litres"\textquotesingle{} is not
assignable to type \textquotesingle"liters"\textquotesingle. \\
11: parameter & Properties are read from one object & Optional fields
for \{\} defaults & TS2345: Argument of type
\textquotesingle\{\}\textquotesingle{} is not assignable to parameter of
type \textquotesingle\{ liters: number; \}\textquotesingle. \\
12: Map loop & Entries contain key and value & Map\textless string,
number\textgreater{} & TS2345: Argument of type
\textquotesingle string\textquotesingle{} is not assignable to parameter
of type \textquotesingle number\textquotesingle. \\
T5: missing property & A wrong key reads undefined in JS & Correct
property spelling & TS2339: Property
\textquotesingle litres\textquotesingle{} does not exist on type
\textquotesingle\{ liters: number; \}\textquotesingle. \\
P7: children & React decides runtime rendering & ReactNode & TS2322:
Type \textquotesingle() =\textgreater{} string\textquotesingle{} is not
assignable to type \textquotesingle ReactNode\textquotesingle. \\
P7: DOM forwarding & Props spread in order &
\texttt{ComponentProps\textless{}\allowbreak{}'button'\textgreater{}\allowbreak{}}
& TS2322: Type \textquotesingle string\textquotesingle{} is not
assignable to type \textquotesingle boolean \textbar{}
undefined\textquotesingle. \\
\end{longtable}
}

\section{Bonus 1: LaTeX}\label{bonus-1-latex}

The complete guide, including all twelve forms, pitfalls, T1-T5, P1-P7,
the comparison table, the depth experiments, the real-code cases and
this log, is included in output/pdf/guide.tex and output/pdf/guide.pdf.
The LaTeX is generated from this README rather than a shortened version.
Code uses listings, the comparison uses a real longtable, and the
contents page links to the sections.

Tools: Pandoc 3.14 and Tectonic 0.17.0. With both executables on PATH,
rebuild from the project root:

\begin{lstlisting}
powershell -ExecutionPolicy Bypass -File scripts/build-guide.ps1
\end{lstlisting}

The script also accepts -Pandoc and -Tectonic executable paths.
scripts/latex-layout.lua sets wrapping paragraph widths for the real
comparison table. It runs Pandoc with -\/-from=gfm -\/-to=latex
-\/-standalone -\/-listings -\/-toc, then Tectonic with -\/-keep-logs.
The build log is output/pdf/guide.log. The output PDF is visually
checked after compilation. Runtime code needs only the npm dependencies;
rebuilding the PDF additionally needs these two document tools and a
first-run TeX package download.

\section{Bonus 2: Depth}\label{bonus-2-depth}

These experiments investigate iteration, evaluation order, copying and
type inference. Predictions were recorded before execution in
experiments.md. They run with npm run all.

\subsection{D1: An array pattern asks for an
iterator}\label{d1-an-array-pattern-asks-for-an-iterator}

\begin{lstlisting}
const [letter, ...letters] = 'mint'
const [crop, ...otherCrops] = new Set(['basil', 'mint', 'basil', 'sage'])
const [[bed, liters]] = new Map([['B2', 2], ['C1', 3]])
console.log('D1 iterables:', letter, letters.join(''), '|', crop, JSON.stringify(otherCrops), '|', bed, liters)
const events = []
function* harvest() {
  try { events.push('start'); yield 'basil'; events.push('second'); yield 'mint' }
  finally { events.push('closed') }
}
const [firstCrop] = harvest()
console.log('D1 early close:', firstCrop, JSON.stringify(events))
try { const [first] = { 0: 'basil', length: 1 }; console.log(first) }
catch (error) { console.log('D1 deliberate experiment:', error.name + ': ' + error.message) }
\end{lstlisting}

\begin{lstlisting}
D1 iterables: m int | basil ["mint","sage"] | B2 2
D1 early close: basil ["start","closed"]
D1 deliberate experiment: TypeError: {(intermediate value)(intermediate value)} is not iterable
\end{lstlisting}

A string yields characters, a Set yields unique entries in insertion
order, and a Map yields key/value pairs. The nested Map pattern takes
the first entry apart. The generator trace contains start and closed but
not second: taking one item closes the unfinished iterator and runs
finally. A numeric property plus length does not make an object
iterable; the deliberate array-like experiment throws. This matters for
generators holding resources: a short pattern can end iteration early.
Array rest would consume the remaining iterator instead.

\subsection{D2: A default is evaluated only when
needed}\label{d2-a-default-is-evaluated-only-when-needed}

\begin{lstlisting}
const defaults = []
function fallback(label, value) { defaults.push(label); return value }
const [pots = fallback('pots', 3), capacity = fallback('capacity', pots * 2),
  explicitNull = fallback('null', 9), zero = fallback('zero', 9)] = [undefined, undefined, null, 0]
console.log('D2 evaluation:', pots, capacity, explicitNull, zero, JSON.stringify(defaults))
\end{lstlisting}

\begin{lstlisting}
D2 evaluation: 3 6 null 0 ["pots","capacity"]
\end{lstlisting}

The pots fallback runs first. The capacity fallback can then read pots
and calculates six. Null and zero are present values, so their fallback
calls never execute. The trace records only pots and capacity. I would
avoid expensive or state-changing defaults unless this order is
intentional; the same pattern can do different work depending on the
incoming values.

\subsection{D3: Object rest copies values, not
descriptors}\label{d3-object-rest-copies-values-not-descriptors}

\begin{lstlisting}
const tag = Symbol('tag')
let reads = 0
const record = Object.create({ inherited: 'not copied' })
Object.defineProperties(record, {
  id: { value: 'B2', enumerable: true },
  crop: { get() { reads++; return 'basil' }, enumerable: true },
  sensor: { value: { moisture: 34 }, enumerable: true },
  hidden: { value: 'not copied', enumerable: false },
  [tag]: { value: 'outdoor', enumerable: true }
})
const { id, ...rest } = record
rest.sensor.moisture = 10
console.log('D3 rest:', id, JSON.stringify(Reflect.ownKeys(rest).map(String)), reads,
  rest[tag], 'hidden' in rest, 'inherited' in rest, record.sensor.moisture)
\end{lstlisting}

\begin{lstlisting}
D3 rest: B2 ["crop","sensor","Symbol(tag)"] 1 outdoor false false 10
\end{lstlisting}

The own enumerable crop, sensor and symbol properties are copied; id is
excluded by the pattern. Inherited and non-enumerable properties are
absent. The crop getter runs once when rest reads its value, and rest
receives a data property rather than the getter itself. JSON alone would
hide symbol keys, so Reflect.ownKeys is used to inspect the result. The
sensor is still shared: mutating the rest sensor changes the original
reading to 10.

\subsection{D4: Destructuring can assign to properties and catch
bindings}\label{d4-destructuring-can-assign-to-properties-and-catch-bindings}

\begin{lstlisting}
const targets = { east: 'basil', west: 'mint' }
;[targets.east, targets.west] = [targets.west, targets.east]
console.log('D4 targets:', JSON.stringify(targets))
try { throw new Error('B2 needs water') }
catch ({ message }) { console.log('D4 catch:', message) }
\end{lstlisting}

\begin{lstlisting}
D4 targets: {"east":"mint","west":"basil"}
D4 catch: B2 needs water
\end{lstlisting}

The right-hand array is evaluated before either target is assigned, so
the two existing object properties swap safely. There are no new
east/west local variables. The catch parameter reads message from the
thrown Error. This is convenient for a known object, but destructuring
an arbitrary thrown null value would itself fail; a defensive TypeScript
catch normally narrows unknown before reading message.

\subsection{D5: Pulling apart types and generic component
props}\label{d5-pulling-apart-types-and-generic-component-props}

\begin{lstlisting}
import type { ReactNode } from 'react'
import { renderToStaticMarkup } from 'react-dom/server'
// Bonus 2: destructuring types, not runtime values.
type First<T extends readonly unknown[]> = T extends readonly [infer H, ...unknown[]] ? H : never
type Rest<T extends readonly unknown[]> = T extends readonly [unknown, ...infer R] ? R : []
function watering(crop: string, liters: number) { return crop + ': ' + liters + ' L' }
type WateringArguments = Parameters<typeof watering>
const head: First<readonly [string, number]> = 'basil'
const tail: Rest<readonly [string, number]> = [2]
const args: WateringArguments = [head, tail[0]]
console.log('D5 types:', head, JSON.stringify(tail), watering(...args))
function CropList<T>({ items, renderItem }: { items: readonly T[]; renderItem: (item: T) => ReactNode }) {
  return <ul>{items.map((item, index) => <li key={index}>{renderItem(item)}</li>)}</ul>
}
console.log('D5 generic props:', renderToStaticMarkup(
  <CropList items={[{ crop: 'basil', liters: 2 }]} renderItem={({ crop, liters }) => <b>{crop}: {liters} L</b>} />
))
\end{lstlisting}

\begin{lstlisting}
D5 types: basil [2] basil: 2 L
D5 generic props: <ul><li><b>basil: 2 L</b></li></ul>
\end{lstlisting}

First extracts the head of a nonempty readonly tuple; Rest infers the
remaining tuple. They operate only during type checking.
First\textless{[}{]}\textgreater{} becomes never and
Rest\textless{[}{]}\textgreater{} becomes {[}{]}, rather than inventing
a runtime value. Parameters produces {[}crop: string, liters: number{]}.
Spreading that tuple supplies the original function arguments. CropList
infers the item object type from items; its renderItem parameter
destructures crop and liters with that inferred type. The generic type
connects the data to the rendering callback without an unsafe cast.

References:
\href{https://tc39.es/ecma262/2024/multipage/ecmascript-language-expressions.html\#sec-destructuring-assignment}{ECMAScript
destructuring assignment and iterator closing},
\href{https://www.typescriptlang.org/docs/handbook/2/conditional-types.html}{TypeScript
conditional types and infer}, and
\href{https://www.typescriptlang.org/docs/handbook/utility-types.html\#parameterstype}{Parameters
utility type}. The iterator and conditional-type explanations above
follow these references; the output is from the local runs.

\section{Bonus 3: Real code}\label{bonus-3-real-code}

These six cases come from the Assignment 3 Task Manager starter, the
installed TypeScript 5.8.3 package, and p-map v7.0.3. Source locations
are recorded explicitly; the examples below are analyses of existing
code. The rewrites show the same extraction behavior within the original
surrounding code, rather than standalone replacements for the whole
application.

\subsection{1. Task Manager state
tuple}\label{1-task-manager-state-tuple}

Source: Downloads/starter.zip -\textgreater{} starter/src/App.tsx, line
37.

Form 1: values by position.

Original extraction:

\begin{lstlisting}
const [tasks, dispatch] = useReducer(taskReducer, [])
\end{lstlisting}

Without destructuring:

\begin{lstlisting}
const statePair = useReducer(taskReducer, [])
const tasks = statePair[0]
const dispatch = statePair[1]
\end{lstlisting}

useReducer returns the current task list and dispatch function. The
pattern names the two positions and keeps the hook call to one
evaluation. This improves readability because React hook pairs have
stable positional roles.

\subsection{2. Task Manager React
props}\label{2-task-manager-react-props}

Source: Downloads/starter.zip -\textgreater{} starter/src/TaskRow.tsx,
line 5.

Forms 5 and 11: object properties in a parameter.

Original extraction:

\begin{lstlisting}
export default function TaskRow({ task, onNext, onDelete }: TaskRowProps) {
\end{lstlisting}

Without destructuring:

\begin{lstlisting}
export default function TaskRow(props: TaskRowProps) {
  const task = props.task
  const onNext = props.onNext
  const onDelete = props.onDelete
  // The original component body follows.
}
\end{lstlisting}

task is data; onNext and onDelete are callback props. There is no
children field in TaskRowProps and the call in App.tsx lines 25-30 is
self-closing. The buttons and task title are created inside TaskRow, so
they are not children supplied by its caller. React key belongs to
element identity and is not an ordinary TaskRow prop.

\subsection{3. TypeScript pattern
formatting}\label{3-typescript-pattern-formatting}

Source: node\_modules/typescript/lib/typescript.js, line 3631
(TypeScript 5.8.3).

Forms 5 and 11: named fields in a parameter.

Original extraction:

\begin{lstlisting}
function patternText({ prefix, suffix }) {
\end{lstlisting}

Without destructuring:

\begin{lstlisting}
function patternText(pattern) {
  const prefix = pattern.prefix
  const suffix = pattern.suffix
  return `${prefix}*${suffix}`
}
\end{lstlisting}

The function only needs prefix and suffix. A short parameter pattern
makes its inputs visible without an extra alias. The plain form makes
the original pattern object available if the function later needs more
properties.

\subsection{4. TypeScript nested watcher
entry}\label{4-typescript-nested-watcher-entry}

Source: node\_modules/typescript/lib/typescript.js, line 7802
(TypeScript 5.8.3).

Form 9: object -\textgreater{} array -\textgreater{} nested object.

Original extraction:

\begin{lstlisting}
const { value: [dirPath, { dirName, options, fileNames }] } = result;
\end{lstlisting}

Without destructuring:

\begin{lstlisting}
const entry = result.value
const dirPath = entry[0]
const details = entry[1]
const dirName = details.dirName
const options = details.options
const fileNames = details.fileNames
\end{lstlisting}

This is the case where I find destructuring less readable. One line
crosses three shapes and hides the intermediate watcher record. The
surrounding code first asserts that the iterator is not done. Naming
entry and details makes that assumption and the record boundary easier
to inspect. The shorter pattern is valid, but the explicit version is
easier to debug when a watcher entry is malformed.

\subsection{5. TypeScript version
parser}\label{5-typescript-version-parser}

Source: node\_modules/typescript/lib/typescript.js, line 4872
(TypeScript 5.8.3).

Forms 2 and 4: skip a regex match and default missing capture groups.

Original extraction:

\begin{lstlisting}
const [, major, minor = "0", patch = "0", prerelease = "", build2 = ""] = match;
\end{lstlisting}

Without destructuring:

\begin{lstlisting}
const major = match[1]
const minor = match[2] === undefined ? '0' : match[2]
const patch = match[3] === undefined ? '0' : match[3]
const prerelease = match[4] === undefined ? '' : match[4]
const build2 = match[5] === undefined ? '' : match[5]
\end{lstlisting}

The first comma skips the complete regex match. Later positions are
capture groups. The defaults handle absent groups and retain explicit
empty strings. This is compact but positional: changing the regular
expression requires checking every capture index.

\subsection{6. p-map dependent option
defaults}\label{6-p-map-dependent-option-defaults}

Source:
\href{https://github.com/sindresorhus/p-map/blob/c9831c7a9b7baa6e986428cc7f1c83af694b5c8e/index.js\#L177}{p-map
v7.0.3 index.js, lines 177-184}.

Forms 7 and 11: defaults in an options parameter.

Original extraction:

\begin{lstlisting}
export function pMapIterable(
    iterable,
    mapper,
    {
        concurrency = Number.POSITIVE_INFINITY,
        backpressure = concurrency,
    } = {},
) {
\end{lstlisting}

Without destructuring:

\begin{lstlisting}
function readOptions(options = {}) {
  const concurrency = options.concurrency === undefined
    ? Number.POSITIVE_INFINITY : options.concurrency
  const backpressure = options.backpressure === undefined
    ? concurrency : options.backpressure
  return { concurrency, backpressure }
}
\end{lstlisting}

In pMapIterable, backpressure defaults to the already-bound concurrency.
Omitting the whole options argument uses \{\}, and omitting either field
uses its default. Null does not activate a default. The later validation
in the actual library checks allowed numeric ranges; destructuring
itself does not validate them. The rewrite isolates option extraction
and leaves that surrounding validation unchanged.

Source notes: the Task Manager cases were read from the original local
starter.zip, not from the completed Assignment 3 repository. Installed
TypeScript source is pinned by package-lock.json; its project is
\href{https://github.com/microsoft/TypeScript}{microsoft/TypeScript}.
The p-map reference is pinned to commit
c9831c7a9b7baa6e986428cc7f1c83af694b5c8e, so future edits cannot shift
this cited version. The small excerpts retain their original source
spelling.

\section{Bonus 4: Git repository}\label{bonus-4-git-repository}

\href{https://github.com/Hurulayn0/assignment4-field-guide}{assignment4-field-guide}

The repository is public: the instructor can read its source and history
without an invitation. node\_modules, local npm caches and disposable
compiler sources are ignored. A clean clone is verified with npm
install, npm run all and npm run check.

The history includes the initial plain JavaScript/TypeScript field guide
on 15 November 2026, followed by separate React implementation, depth
experiments, real-code analysis and final LaTeX/package verification
stages on 16 November 2026. These are meaningful completed stages on
their actual dates; no timestamps are backdated. The final published
history contains at least five commits across those two dates. The
instructor decides the bonus score.

\section{My own log}\label{my-own-log}

This assignment took a lot of my time. The points needing the most care
in this part were defaults, null values, tuple positions and possibly
missing array items.

The entries below record actual experiment results. Predictions were
recorded before running the checks.

\begin{enumerate}
\def\labelenumi{\arabic{enumi}.}
\item
  Missing values and undefined used the defaults in forms 4 and 7, while
  null stayed null. A default does not replace every empty-looking
  value.
\item
  Deliberate experiment: destructuring undefined and null threw
  TypeError. Both errors were caught; the exact messages are in the
  pitfalls section above and output.txt.
\item
  Deliberate experiment: a bare object assignment produced
  \texttt{SyntaxError:\allowbreak{} Unexpected token '=\allowbreak{}'}.
  Parentheses fixed the parse, and the corrected result was 2.
\item
  The sensor binding remained shared. Changing its moisture value
  changed the original record too; rest creates only a shallow outer
  copy.
\item
  The empty typed array gave undefined. The default produced empty, and
  the undefined check avoided calling toUpperCase on a missing value.
\item
  Deliberate T5 experiment: the misspelled property produced TS2339. The
  wrong key gave undefined at runtime and the corrected key gave 2. The
  exact compiler diagnostic is in logs/t5-diagnostic.txt.
\item
  In T1, the variable named string contained basil. The colon inside the
  pattern renamed the binding; the type annotation belongs after the
  whole pattern.
\item
  A direct function child produced the exact React warning in P3 and no
  visible content. The render-prop component called its function and
  rendered mint. The deliberately invalid object child threw a caught
  Error.
\item
  All five P7 misuse cases failed strict checks as predicted. Exact
  messages are recorded in logs/compiler-diagnostics.txt.
\item
  The generator closed after one value and ran finally, while the
  array-like object was rejected. Object rest copied a symbol key that
  JSON.stringify would not reveal; Reflect.ownKeys made it visible. Only
  the two undefined positions called the D2 fallback.
\end{enumerate}

\end{document}
